% !Mode:: "TeX:UTF-8"
\chapter{交通流量预测模型实现与结果分析}

在本章中，将详细介绍本研究的核心预测算法——PDFormer (Propagating Diffusion Transformer)。通过在构建的交通数据集上进行模型训练与参数调优，本章旨在验证所设计的时空预测引擎在复杂多模态交通场景下的准确性与鲁棒性。

\section{PDFormer 模型概述}
PDFormer（Propagating Diffusion Transformer）作为本研究的核心时序预测引擎，其设计初衷是为了解决传统时空模型在处理长程依赖与动态空间扩散时的局限性。该模型在架构上深度融合了图神经网络（GNN）的拓扑感知能力与 Transformer 架构的全局建模优势，构建了一套完备的时空双重对齐逻辑。在空间维度上，模型通过引入多头空间自注意力机制（Spatial Self-Attention），能够自适应地捕获路网中非相邻传感器间的隐式拓扑关联，这弥补了传统 GCN 仅能提取局部邻接特征的不足。在时间维度上，PDFormer 利用时间自注意力机制对交通流的日周期与周周期规律进行深度挖掘，确保了模型在长时程预测任务中依然具备极高的特征保持度。

针对真实城市交通中普遍存在的“拥堵滞后扩散”现象，PDFormer 创新性地引入了传播扩散机制。该机制通过扩散卷积逻辑模拟了交通压力在物理路网中的传递过程，能够精确刻画某一节点流量剧增对邻近乃至远端路口产生的时空连锁反应。为了进一步强化对路网全局结构的感知，模型还集成了基于拉普拉斯矩阵的特征嵌入技术（Laplacian PE）。通过对路网拉普拉斯矩阵进行谱分解，提取特征向量作为空间位置编码，模型实现了对交通网络微观动态特征与宏观拓扑拓扑结构的同步表征。这种多维度的特征集成体系，为 PDFormer 在面对突发事件导致的非线性交通波动时提供了极强的稳健性。

通过在国际公认的交通基准数据集 PeMS04、PeMS07 及 PeMS08 上的实验结果（如表 5-1 所示）可以看出，PDFormer 相比于传统的向量自回归模型（VAR）、支持向量机（SVR）以及先进的时空图卷积网络（如 STGCN、GWNET、STGNCDE 等）进行综合评估，可以清晰地观察到该模型在平均绝对误差（MAE）、平均绝对百分比误差（MAPE）及均方根误差（RMSE）等核心指标上均取得了最优表现（SOTA）。特别是在路网结构最为复杂且采样点最为密集的 PeMS07 数据集上，PDFormer 的预测误差显著低于其他对比模型，这充分证明了其在处理大规模、异构时空特征时的卓越性能。


\begin{table}[htbp]
  \centering
  \caption{Performance on Graph-based Datasets.}
  \label{tab:performance_graph}
  \resizebox{0.8\textwidth}{!}{
  \begin{tabular}{l|ccc|ccc|ccc}
    \toprule
    \multirow{2}{*}{Model} & \multicolumn{3}{c|}{PeMS04} & \multicolumn{3}{c|}{PeMS07} & \multicolumn{3}{c}{PeMS08} \\
    \cline{2-10}
    & MAE & MAPE(\%) & RMSE & MAE & MAPE(\%) & RMSE & MAE & MAPE(\%) & RMSE \\
    \midrule
    VAR     & 23.750 & 18.090 & 36.660 & 101.200 & 39.690 & 155.140 & 22.320 & 14.470 & 33.830 \\
    SVR     & 28.660 & 19.150 & 44.590 & 32.970  & 15.430 & 50.150  & 23.250 & 14.710 & 36.150 \\
    \midrule
    DCRNN   & 22.737 & 14.751 & 36.575 & 23.634  & 12.281 & 36.514  & 18.185 & 11.235 & 28.176 \\
    STGCN   & 21.758 & 13.874 & 34.769 & 22.898  & 11.983 & 35.440  & 17.838 & 11.211 & 27.122 \\
    GWNET   & 19.358 & 13.301 & 31.719 & 21.221  & 9.075  & 34.117  & 15.063 & 9.514  & 24.855 \\
    MTGNN   & 19.076 & 12.961 & 31.564 & 20.824  & 9.032  & 34.087  & 15.396 & 10.170 & 24.934 \\
    STSGCN  & 21.185 & 13.882 & 33.649 & 24.264  & 10.204 & 39.034  & 17.133 & 10.961 & 26.785 \\
    STFGNN  & 19.830 & 13.021 & 31.870 & 22.072  & 9.212  & 35.805  & 16.636 & 10.547 & 26.206 \\
    STGODE  & 20.849 & 13.781 & 32.825 & 22.976  & 10.142 & 36.190  & 16.819 & 10.623 & 26.240 \\
    STGNCDE & 19.211 & 12.772 & 31.088 & 20.620  & 8.864  & 34.036  & 15.455 & 9.921  & 24.813 \\
    \midrule
    STTN    & 19.478 & 13.631 & 31.910 & 21.344  & 9.932  & 34.588  & 15.482 & 10.341 & 24.965 \\
    GMAN    & 19.139 & 13.192 & 31.601 & 20.967  & 9.052  & 34.097  & 15.307 & 10.134 & 24.915 \\
    TFormer & 18.916 & 12.711 & 31.349 & 20.754  & 8.972  & 34.062  & 15.192 & 9.925  & 24.883 \\
    ASTGNN  & 18.601 & 12.630 & 31.028 & 20.616  & 8.861  & 34.017  & 14.974 & 9.489  & 24.710 \\
    \midrule
    \textbf{PDFormer} & \textbf{18.321} & \textbf{12.103} & \textbf{29.965} & \textbf{19.832} & \textbf{8.529} & \textbf{32.870} & \textbf{13.583} & \textbf{9.046} & \textbf{23.505} \\
    \bottomrule
  \end{tabular}
  }
\end{table}

\vspace{2em}

为了验证模型的可迁移性与复现性，本研究在本地环境下针对 PeMS04 和 PeMS08 数据集进行了完整的训练与测试流程。实验结果表明，复现得到的 $MAE$ 与 $RMSE$ 指标与官方基准高度吻合，进一步夯实了本系统采用 PDFormer 作为底座预测引擎的可靠性。

\begin{figure}[htbp]
    \centering
    \subfigure[PeMS04 预测对比图]{
        \includegraphics[width=0.48\textwidth]{image5.png}
    }
    \subfigure[PeMS08 预测对比图]{
        \includegraphics[width=0.48\textwidth]{image6.png}
    }
    \caption{测试集预测值与真实值对比分析}
    \label{fig:predict_compare}
\end{figure}

此外，本研究绘制了模型在测试集上的预测值与真实值对比曲线图（如图 6-1 所示）。两个模型均选用第一个节点在测试机上前三天的预测表现。通过对曲线的观察可以发现，PDFormer 在早晚高峰等流量剧烈波动的时段，依然能够保持极高的拟合精度，准确捕捉到了交通流的非线性变化趋势。

\section{多步预测性能评价与局部验证}
在完成模型的基础复现后，本研究进一步针对 PeMS04 和 PeMS08 数据集进行了精细化的多步预测性能评估。如图 5-X 和 图 5-Y 所示，实验记录了模型在未来 12 个预测步长（对应聚合后的 60 分钟时长）内的误差变化趋势。

\begin{figure}[htbp]
    \centering
    \subfigure[PeMS04 性能评估]{
        \includegraphics[width=0.48\textwidth]{image7.png}
    }
    \subfigure[PeMS08 性能评估]{
        \includegraphics[width=0.48\textwidth]{image8.png}
    }
    \caption{测试集预测值与真实值对比分析}
    \label{fig:performance_evaluation}
\end{figure}

在 PeMS04 数据集的验证中，PDFormer 表现出了极强的预测稳定性。在初始预测步（第 1 步，即未来 5 分钟）时，其平均绝对误差（$MAE$）仅为 16.52，均方根误差（$RMSE$）为 27.03。随着预测时间跨度的增加，由于时间不确定性的累积效应，误差呈现出缓慢上升的趋势。至第 12 个预测步（未来 60 分钟）时，$MAE$ 增至 18.31，$RMSE$ 增至 30.35。这一结果不仅与表 5-1 中的 SOTA 基准高度一致，更揭示了 PDFormer 在处理大规模路网流量时，能够有效通过传播扩散机制抑制远期误差的剧烈发散。

与之类似，在 PeMS08 数据集的测试中，模型展现了更为优异的拟合能力。第 1 步的 $MAE$ 低至 11.76，而在完成一小时（第 12 步）的连续预测后，$MAE$ 仅增长至 13.59。通过对这两组多步预测数据的综合对比可以发现，PDFormer 在前 3 个步长内的精度优势尤为明显，这对于交通智能体处理突发性的路况查询（如“未来 15 分钟的交通演变”）具有极高的实战价值。此外，较低的 $masked\_MAPE$（在 PeMS08 上维持在 8\% 至 9\% 之间）进一步证明了模型在过滤背景噪声、捕捉核心交通流特征方面的卓越性能。这种基于本地环境的详尽复现工作，不仅定量验证了 PDFormer 架构的可靠性，也为后续在本研究实现高精度流量预测提供了坚实的实验证据。


\section{实验配置与训练策略}
为了确保实验结果的科学性与可复现性，本研究构建了统一的硬件与软件实验平台。硬件环境方面，所有实验均在配置有 NVIDIA GeForce RTX 3090 系列显卡（显存 24GB）的服务器上完成，算力支撑足以满足 PDFormer 模型大规模参数迭代的需求。软件环境则基于 Ubuntu 22.04 操作系统，采用 Python 3.9 作为核心开发语言，并依托 LibCity 这一权威的时空数据挖掘框架进行模型训练与评估。

针对不同性质的实验任务，本研究制定了差异化的超参数配置策略。在进行基准性能验证时，针对 PeMS04 和 PeMS08 开源数据集，模型输入维度设定为交通领域标准的单维度；模型选用了 AdamW 作为核心优化器，并配合 10 个 Epoch 的预热（Warmup）阶段，将批处理大小（BatchSize）设定为 16。此外，针对不同路网拓扑的特性，本研究在具体参数上进行了微调。如表 6-2 所示，PeMS04 数据集的聚类分析参考天数设定为 14 天，而 PeMS08 则设定为 21 天，以匹配不同路网下交通流的长周期特征。在时间窗口配置上，实验统一遵循“过去 1 小时观测预测未来 1 小时态势”的逻辑，即将输入与输出步长均设为 12。


\begin{table}[htbp]
  \centering
  \caption{PeMS04 与 PeMS08 实验超参数配置汇总}
  \label{tab:hyperparameters_pems_final}
  \renewcommand{\arraystretch}{1.2}
  \footnotesize
  \begin{tabular}{lcc}
    \toprule
    \textbf{参数项} & \textbf{PeMS04 配置值} & \textbf{PeMS08 配置值} \\
    \midrule
    时间窗口 (Input/Output) & 12 / 12 & 12 / 12 \\
    训练/验证/测试比例 & 6:2:2 & 6:2:2 \\
    批处理大小 (Batch Size) & 16 & 16 \\
    最大迭代轮数 (Max Epoch) & 200 & 200 \\
    优化器类型 & AdamW & AdamW \\
    学习率调度策略 & CosineLR (Warmup: 10) & CosineLR (Warmup: 10) \\
    聚类方法与周期 (Cluster) & kshape (14 days) & kshape (21 days) \\
    \bottomrule
  \end{tabular}
\end{table}

在具体的模型运行流程中，本文通过脚本化指令 run\_model.py 触发任务，并针对不同任务指定相应的配置文件。具体的执行命令如下：

\begin{lstlisting}[
    language=bash,
    caption={启动模型训练命令},
    label={code-vllm-server},
    breaklines=true,
    frame=single,
    numbers=left,
    escapeinside={@}{@}
]
nohup python run_model.py --task traffic_state_pred --model PDFormer
--dataset PeMS04 --config_file PeMS04 > train.log 2>&1 &
\end{lstlisting}


\section{训练过程与收敛性分析}
本研究通过监测训练损失（Training Loss）与验证损失（Evaluation Loss）在迭代过程中的动态演化，对 PDFormer 模型的收敛性能进行了深度评估。如图 6-3 所示，在 PeMS04 数据集的训练过程中，模型展现出了极佳的收敛特征。在训练初期（0 至 20k 个迭代步长），由于 PDFormer 能够迅速识别路网中的基础时空模式，训练损失表现出显著的下降态势。随着迭代次数突破 60k 步，Training Loss 曲线逐渐由快速衰减转为平稳微调，最终稳定在 31.5 左右。

在验证集表现方面，Evaluation Loss 的演化趋势与训练集保持了高度的逻辑一致性。如图 6-3 所示，验证损失在 20k 步左右即实现了从高位（约 200）到低位（约 40）的跨越式下降，并最终平稳收敛于 35 附近。这种“快速收敛、长程稳定”的曲线特征表明，PDFormer 架构中集成的空间自注意力机制与 Laplacian 特征嵌入能够有效加速对复杂路网拓扑的特征提取。

此外，通过对比 Training Loss 与 Evaluation Loss 的最终稳态值可以发现，两条曲线之间维持着极小的泛化间隙（Generalization Gap），且验证集损失未出现任何上扬迹象。这有力地证明了模型在长达 120k 步的深度迭代中，始终保持着优异的泛化能力，并未产生过拟合现象。这种稳健的收敛性能确保了模型在处理非线性、高维度的交通流预测任务时，能够提供具有高度置信度的数值输出，为上层智能体系统的复杂决策提供了稳固的算法基石。

\begin{figure}[htbp]
    \centering
    \subfigure[eval loss 曲线]{
        \includegraphics[width=0.48\textwidth]{image9.png}
    }
    \subfigure[train loss 曲线]{
        \includegraphics[width=0.48\textwidth]{image10.png}
    }
    \caption{PDFormer 模型训练与验证损失曲线分析}
    \label{fig:performance_loss}
\end{figure}

\section{实验结果多维指标评价}
模型最终在测试集上展现了卓越的预测性能。本研究通过平均绝对误差（$MAE$）、平均绝对百分比误差（$MAPE$）以及均方根误差（$RMSE$）三个核心维度对预测精度进行了定量评价。针对 PeMS04 与 PeMS08 数据集的复现实验结果显示，PDFormer 在多步预测任务中表现极其稳健。以 PeMS08 为例，模型在未来 12 个预测步长（共 60 分钟）内的平均 $MAE$ 稳定在 13.59 附近，且起始步长的预测精度极高。值得注意的是，在不同数据集的横向对比中，PDFormer 表现出的 $MAPE$ 始终维持在较低的百分比区间，这证明了模型不仅在数值绝对值上预测精准，在预测交通波动的相对比例方面同样具有显著优势。结合前文所述的 Loss 收敛情况，可以得出结论：PDFormer 通过空间自注意力与传播扩散机制，成功克服了传统模型在长程预测中误差累积较快的问题。这一高精度的数值预测结果，直接支撑了第四章中交通智能体对态势感知的“深度分析”功能，使得系统生成的分析报告具备了极高的工程实用价值与学术参考意义。